\documentclass[titlepage,a4paper]{jsarticle}
\usepackage{../../../sty/import}% 各種パッケージインポート
\usepackage{../../../sty/title}% タイトルページの変更

%% タイトルページの変数
% レポートタイトル
\title{原子力材料と核燃料}
% 提出日
\expdate{\today}
% 科目名
\subject{原子力材料と核燃料}
% 分野
\class{情報経営システム工学分野}
% 学年
\grade{M1}
% 学籍番号
\mynumber{24336488}
% 記述者
\author{本間三暉}

\begin{document}
% titleページ作成
\maketitle
\section{原子力発電の将来について思うことを自由に述べよ}
私は，将来の日本においても原子力発電は重要な役割を果たすべきであると考える．
近年はAI技術が急速に発展しており，将来的にはAIを支えるデータセンターの増加などによって電力需要がさらに高まる可能性がある．
そのような状況において，安定して大量の電力を供給できる電源を確保することは重要であり，原子力発電もその有力な選択肢の一つであると思う．

一方で，福島第一原子力発電所事故のような重大事故を二度と起こさないため，安全対策はこれまで以上に徹底する必要がある．
しかし，原子力発電所は運転を停止していても燃料や設備の管理が必要であり，一定のコストやリスクが存在する．
そのため，安全性が十分に確認された原子力発電所については，再稼働を進めることも重要であると考える．

また，日本は火力発電に必要な燃料の多くを海外からの輸入に依存している．
近年は国際情勢の変化によって燃料価格が大きく変動しており，電気料金にも影響を与えている．
エネルギー安全保障や電気料金の安定化の観点からも，原子力発電を適切に活用していくことが望ましいと考える．

今後は再生可能エネルギーの導入もさらに進むと考えられるが，発電量が天候に左右されることに加え，一部では太陽光発電設備の設置に伴う森林開発や景観への影響なども課題として指摘されている．
そのため，再生可能エネルギーだけに依存するのではなく，原子力発電を含む複数の電源を組み合わせながら，安定した電力供給と環境負荷の低減を両立していくことが重要であると考える．
福島第一原子力発電所事故の教訓を活かしながら安全性をさらに向上させ，将来の電力需要の増加にも対応できる原子力発電のあり方を模索していくべきだと思う．

% 参考文献
\begin{thebibliography}{99}
\bibitem{原子力材料と核燃料 講義資料}原子力材料と核燃料 講義資料
\end{thebibliography}

\end{document}